\begin{ouverture}[raisonner]{12}{Droites remarquables et raisonnement}{Partie 4 · Géométrie et espace}
{Médiatrices · Hauteurs · Médianes · Réciproque · Contraposée}
\objectif{Reconnaître les droites remarquables d’un triangle et utiliser leurs propriétés.}
\objectif{Distinguer hypothèse, conclusion, réciproque et contraposée.}
\objectif{Démontrer par déduction, par un contre-exemple ou par l’absurde.}
\prerequis{Médiatrice, hauteurs, médianes, angles et parallélogrammes (Cinquième).}
\end{ouverture}
\pagestyle{otMain}\essentiel\fontsize{11}{15}\selectfont
\phantomsection\label{cours-12-1}
\essSection{1}{Les droites remarquables}
\begin{definition}Dans un triangle, la médiatrice d’un côté est la droite perpendiculaire à ce côté en son milieu. Une hauteur passe par un sommet et est perpendiculaire à la droite qui porte le côté opposé. Une médiane passe par un sommet et par le milieu du côté opposé. La bissectrice d’un angle le partage en deux angles de même mesure.\end{definition}
\label{fig-cours-12-1-cours-droites-remarquables}\figureGeo{cours-droites-remarquables}

\begin{propriete}[Rappels de Cinquième]Un point appartient à la médiatrice de $[AB]$ si et seulement s’il est équidistant de A et de B. Les trois médiatrices d’un triangle sont concourantes au centre de son cercle circonscrit ; les trois hauteurs sont concourantes, ainsi que les trois médianes. Une médiane partage le triangle en deux triangles de même aire.\end{propriete}
\begin{exemple}Dans un triangle ABC isocèle en A, la médiatrice de $[BC]$ passe par A, car $AB=AC$. Elle porte donc la médiane et la hauteur issues de A, et c’est la bissectrice de l’angle en A : elle est l’axe de symétrie du triangle.\end{exemple}
\clearpage
\phantomsection\label{cours-12-2}
\essSection{2}{Démontrer et raisonner}
\begin{definition}Un énoncé « Si P, alors Q » a une hypothèse P et une conclusion Q. Sa réciproque est l’énoncé « Si Q, alors P ». Sa contraposée est l’énoncé « Si Q est faux, alors P est faux ». Un énoncé démontré s’appelle un théorème.\end{definition}
\begin{propriete}Si un énoncé est vrai, sa contraposée est vraie aussi ; sa réciproque peut être vraie ou fausse. Lorsqu’un énoncé et sa réciproque sont vrais, on obtient une propriété caractéristique : « M est sur la médiatrice de $[AB]$ si et seulement si $MA=MB$ ».\end{propriete}
\begin{exemple}« Si ABCD est un rectangle, alors ses diagonales ont la même longueur » est vrai. Sa réciproque est fausse : un trapèze isocèle a des diagonales de même longueur. Sa contraposée est vraie : si les diagonales ont des longueurs différentes, ABCD n’est pas un rectangle.\end{exemple}
\begin{savoirfaire}[Trois façons de raisonner]Déduire : enchaîner les données, une propriété citée et la conclusion. Réfuter : un seul contre-exemple montre qu’un énoncé général est faux. Raisonner par l’absurde : supposer le contraire de ce que l’on veut prouver et aboutir à une contradiction.\end{savoirfaire}
\begin{attention}Une mesure sur un dessin, ou quelques exemples qui marchent, ne prouvent pas une propriété générale. On précise si une propriété utilisée est admise ou démontrée.\end{attention}

\clearpage\rubrique{Découverte}
\begin{blocLivre}{ex-12-1}\begin{exercice}[1]{Reconnaître des droites}
Sur la figure codée, nommer la nature des droites $(d_1)$, $(d_2)$ et $(d_3)$ dans le triangle ABC, en justifiant par le codage.
\label{fig-ex-12-1}\figureGeo{rem-01}
\end{exercice}\end{blocLivre}
\begin{blocLivre}{ex-12-2}\begin{exercice}[1]{Hypothèse et conclusion}
Énoncé : « Si ABCD est un rectangle, alors ses diagonales ont la même longueur. » Donner l’hypothèse et la conclusion, puis écrire la réciproque. Est-elle vraie ?
\label{fig-ex-12-2}\figureGeo{rem-02}
\end{exercice}\end{blocLivre}
\begin{blocLivre}{ex-12-3}\begin{exercice}[1]{Une hauteur et une aire}
Dans le triangle ABC, $AB=8$ cm et la hauteur issue de C mesure $CH=3$ cm. Calculer l’aire du triangle. Pourquoi ne peut-on pas utiliser AC à la place de CH ?
\label{fig-ex-12-3}\figureGeo{rem-03}
\end{exercice}\end{blocLivre}
\begin{blocLivre}{ex-12-4}\begin{exercice}[1]{Un point de la médiatrice}
Le point M appartient à la médiatrice de $[AB]$ et $MA=4{,}5$ cm. Que vaut MB ? Un point N vérifie $NA=NB=6$ cm : à quelle droite appartient-il ?
\label{fig-ex-12-4}\figureGeo{rem-04}
\end{exercice}\end{blocLivre}
\begin{blocLivre}{ex-12-5}\begin{exercice}[1]{Les angles d’un triangle isocèle}
ABC est isocèle en A, avec $AB=AC$ et $\widehat{BAC}=74^\circ$. Déterminer les deux autres angles.
\label{fig-ex-12-5}\figureGeo{rem-05}
\end{exercice}\end{blocLivre}

\rubrique{Entraînement}
\begin{blocLivre}{ex-12-6}\begin{exercice}[2]{Un contre-exemple}
Léo affirme : « Deux triangles qui ont la même aire ont les mêmes longueurs de côtés. » Montrer qu’il a tort.
\label{fig-ex-12-6}\figureGeo{rem-06}
\end{exercice}\end{blocLivre}
\begin{blocLivre}{ex-12-7}\begin{exercice}[2]{Une médiane partage l’aire}
Dans le triangle ABC, I est le milieu de $[BC]$ et H le pied de la hauteur issue de A. L’aire de ABC vaut 30 cm$^2$. Démontrer que ABI et AIC ont la même aire, puis la calculer.
\label{fig-ex-12-7}\figureGeo{rem-07}
\end{exercice}\end{blocLivre}
\begin{blocLivre}{ex-12-8}\begin{exercice}[2]{Utiliser la contraposée}
On sait que : « Si un point appartient à la médiatrice de $[AB]$, alors il est équidistant de A et de B. » On mesure $PA=5$ cm et $PB=6$ cm. Que peut-on conclure pour P ? Quel énoncé utilise-t-on ?
\label{fig-ex-12-8}\figureGeo{rem-08}
\end{exercice}\end{blocLivre}
\begin{blocLivre}{ex-12-9}\begin{exercice}[2]{Une réciproque fausse}
L’énoncé « Si un quadrilatère est un losange, alors ses diagonales sont perpendiculaires » est vrai. Écrire sa réciproque et montrer qu’elle est fausse.
\label{fig-ex-12-9}\figureGeo{rem-09}
\end{exercice}\end{blocLivre}
\begin{blocLivre}{ex-12-10}\begin{exercice}[2]{Le triangle isocèle}
ABC est isocèle en A et I est le milieu de $[BC]$. Démontrer que A appartient à la médiatrice de $[BC]$. En déduire que $(AI)$ est à la fois la médiane et la hauteur issues de A.
\label{fig-ex-12-10}\figureGeo{rem-10}
\end{exercice}\end{blocLivre}

\rubrique{Problèmes}
\begin{blocLivre}{ex-12-11}\begin{exercice}[3]{Le point commun des hauteurs}
Dans le triangle ABC, les hauteurs issues de A et de B se coupent en H. Démontrer que $(CH)$ est perpendiculaire à $(AB)$. Si ABC est rectangle en C, où se trouve H ?
\label{fig-ex-12-11}\figureGeo{rem-11}
\end{exercice}\end{blocLivre}
\begin{blocLivre}{ex-12-12}\begin{exercice}[3]{Raisonner par l’absurde}
Démontrer par l’absurde qu’un triangle ne peut pas avoir deux angles droits.
\label{fig-ex-12-12}\figureGeo{rem-12}
\end{exercice}\end{blocLivre}
\begin{blocLivre}{ex-12-13}\begin{exercice}[3]{Trois maisons et un puits}
Trois maisons A, B et C ne sont pas alignées. On veut creuser un puits à égale distance des trois maisons. Expliquer comment construire son emplacement O et justifier qu’il convient.
\label{fig-ex-12-13}\figureGeo{rem-13}
\end{exercice}\end{blocLivre}
\begin{blocLivre}{ex-12-14}\begin{exercice}[3]{Vrai ou faux ?}
Dire si chaque énoncé est vrai ou faux et justifier.\begin{questions}
\item Si deux droites sont perpendiculaires à une même droite, alors elles sont parallèles.
\item Si un triangle a deux côtés de même longueur, alors ses trois angles sont égaux.
\item Si deux droites ne sont pas parallèles, alors elles ne sont pas perpendiculaires à une même droite.
\end{questions}
\label{fig-ex-12-14}\figureGeo{rem-14}
\end{exercice}\end{blocLivre}
\begin{blocLivre}{ex-12-15}\begin{exercice}[3]{Une médiane dans une médiane}
Dans le triangle ABC, d’aire 36 cm$^2$, I est le milieu de $[BC]$ et J le milieu de $[AI]$. Démontrer que l’aire du triangle ABJ est le quart de celle de ABC, puis la calculer.
\label{fig-ex-12-15}\figureGeo{rem-15}
\end{exercice}\end{blocLivre}

\clearpage
